% Track R Supplement S3, revised theory-owned integration draft (2026-10-30).
% Source: earlier TH_SUPPL_S3_SELECTION_BOUNDS_DROPIN_20261029.tex,
% A8 calibration/theory, registered v1/v2/v2-E1/v3 files and released
% FAST64, ImageNet D1-P2, SITCOM B16 and DAPS F8 comparison relays.
% Do NOT re-fit a registered prediction or read xp/returns for this draft.
% Copy into paper only after planner approves the integration and labels.
\section{Residual selection, exact finite-panel evaluation, and the identification limit}
\label{sec:suppl-selection}

\paragraph{Candidate classes and the exchangeable-rank reference.}
For a fixed image and solver, label a generated candidate by
$J\in\{G,O_g,O_p,C,I\}$, denoting raw GOOD, globally rotated
ORBIT, partially rotated ORBIT, measurement-CONSISTENT but wrong,
and measurement-INCONSISTENT, with class probabilities
$\theta=(p,o_g,o_p,c,u)$.
Set $q=p+o_g+o_p$ and $s=q+c=1-u$.
The conventional exchangeable-passing-ranks model (S2) takes all
non-$I$ candidates to have equal chance to win the physical-residual
ranking, while $I$ never wins when a passing candidate exists (S1).
Let $h$ be the chance that the orientation verifier breaks GOOD and
$\phi_g,\phi_p$ the repair chances for $O_g,O_p$, respectively. With
$a=(1-h)p+\phi_go_g+\phi_po_p$, independent fresh-$K$ sampling gives
\begin{align}
 E_0^{\mathrm{S2}}(K)&=(1-u^K)\frac{p}{s}, &
 E_1^{\mathrm{S2}}(K)&=(1-u^K)\frac{a}{s},\label{eq:suppl-S3-S2}\\
 O_{\mathrm{raw}}(K)&=1-(1-p)^K, &
 O_{\mathrm{orbit}}(K)&=1-(1-q)^K.\label{eq:suppl-S3-oracle}
\end{align}
The formulas use a fresh-$K$ iid model and assume $s>0$;
for $s=0$, $E_0=E_1=0$ under the stated $I$-last rule.
They are neither the physical law of native diffusion residuals nor
the exact evaluation of an existing fixed panel.

\paragraph{A class-conditional rank race and exact winner law.}
For a residual-\emph{rank} surrogate, draw independent competing
clocks $T\mid J=j\sim\operatorname{Exp}(w_j)$ with
$w_G=1$, $w_{O_g}=w_{O_p}=w_O$, $w_C>0$ and $w_I=0$.
A smaller clock is treated as a lower residual. For a class with
$w_j>0$, the probability it wins among $K$ fresh candidates is
\begin{equation}\label{eq:suppl-race-integral}
 W_j(K\mid\theta,w)
 =K\theta_j\int_0^\infty w_j e^{-w_jt}
 \left(\sum_\ell\theta_\ell e^{-w_\ell t}\right)^{K-1}\,dt.
\end{equation}
To prove this, designate one of the $K$ runs as the winning class-$j$
run: its joint class/time density is $\theta_j w_j e^{-w_jt}$,
while each other run has survival function
$\sum_\ell\theta_\ell e^{-w_\ell t}$.
For realized class counts $N_j$, competing exponential minima
instead give the elementary Plackett--Luce identity
\begin{equation}\label{eq:suppl-race-conditional}
 \Pr(J_{\arg\min T}=j\mid (N_\ell)_\ell)
 =\frac{N_jw_j}{\sum_\ell N_\ell w_\ell}
\end{equation}
when the denominator is positive. If every clock rate is zero, a
predeclared deterministic tie convention is required. Thus, under
class-conditional independent repair decisions,
\begin{equation}\label{eq:suppl-race-policies}
 E_0^{\mathrm{race}}(K)=W_G(K),\qquad
 E_1^{\mathrm{race}}(K)=(1-h)W_G(K)
 +\phi_gW_{O_g}(K)+\phi_pW_{O_p}(K).
\end{equation}
The probabilities refer to a single chosen candidate, \emph{not}
repairing the entire pool or using ground truth for selection.
Taking $K\to\infty$ at fixed class rates yields
\[
 E_0^{\mathrm{race}}(\infty)
 =\frac{p w_G}{p w_G+o_gw_O+o_pw_O+cw_C}
\]
if the denominator is positive, generally different from $p/s$.
Exponential-race formulas are classical; their role here is a
solver-specific calibrated approximation to physical-residual ranks.
An exact Fourier twin has the same measurement magnitude, but two
\emph{separately generated} G and O outputs can have different
non-orbit reconstruction errors and therefore different residuals.

\paragraph{What is and is not learned from a few ranks.}
The authorized A8 calibration comprises 540 eligible image--solver
rows and 1,960 candidates across DAPS, SITCOM, and FAST-DIPS:
1,171 G, 58 global O, one partial O, 215 C and 515 I.
Within-image G ranks ahead of C in 200/205 available pair comparisons;
E0 chooses G in 92/100 image--solver rows with mixed passing classes,
versus an expectation of 51.92 under the exchangeable S2 reference.
These pair comparisons share runs/images and are \emph{not}
independent binomial trials. Composite Bradley--Terry fits regularize
$\log(w_O),\log(w_C)$ independently with a zero-mean Gaussian
logweight penalty (primary scale 2.5); they tie P to O because only
one partial-orbit candidate is observed. The approximate primary
class rates $(w_O,w_C)$ are DAPS $(.1704,.0201)$, SITCOM
$(.9536,.0454)$, and FFHQ FAST $(.0366,.0236)$.
An $I$-last surrogate has observed exceptions, handled in the
frozen scripts by a small finite clock rate only for the affected
\emph{observed} prefix candidate. Do not claim those fitted weights
are physically calibrated units of a Fourier residual or generally
portable to other datasets.

The oracle--selector gap need \emph{not} be monotone in $K$ under
class-dependent ranks, even when the GOOD rate is larger.
With $p_G=p_{\mathrm{bad}}=1/2$ and
$(w_G,w_{\mathrm{bad}})=(1,.03)$,
\[
 D_K=\sum_{g=1}^{K-1}\binom Kg2^{-K}
 \frac{(K-g)\,.03}{g+(K-g)\,.03}
\]
has $D_6=.0386195271>D_7=.0385396532$ and
$D_\infty=.0291262136$. The monotone-gap observation for the
exchangeable reference must therefore \emph{not} be transferred
into a general race theorem.

\paragraph{Exact all-subset policy and oracle estimands.}
For a \emph{completed fixed panel} of $n$ candidates ordered by the
measured residual, let rank $r=1$ be the best. If a uniform $K$-subset
is drawn, the chance that its minimum-residual candidate has rank
$r$ is exactly
\begin{equation}\label{eq:suppl-subset-weight}
 \omega_r(n,K)=\frac{\binom{n-r}{K-1}}{\binom nK},
\end{equation}
with the binomial coefficient defined as zero when its arguments are
infeasible. For observed per-run success indicators $b_r$ after
any specified orientation action, the \emph{complete U-statistic}
policy success is $\sum_{r=1}^n\omega_r(n,K)b_r$.
If precisely $g$ of $n$ candidates are accepted by a ground-truth
oracle, its complete U-statistic is
\begin{equation}\label{eq:suppl-oracle-subset}
 \widehat O(K)=1-\frac{\binom{n-g}{K}}{\binom nK}.
\end{equation}
These identities use no statistical extrapolation and are valid for
the \emph{actual completed panel}. They prevent conflating a fixed
panel with independently regenerated pools for each $K$.

\paragraph{Exact conditional-prefix oracle means.}
For a registered conditional-prefix forecast, suppose $n_0$ old
runs with fixed class counts $C$ remain in a panel of size $n$;
only $m=n-n_0$ future runs are generated. Let
$\theta\mid C\sim\operatorname{Dirichlet}(C+\alpha)$ and let the
oracle's target be either G alone or $\{G,O_g,O_p\}$.
Write $g_0$ for old successes, $A$ for the posterior total
concentration, and $A_{\neg}$ for its non-target concentration.
Then the \emph{exact model-conditional} expected oracle, without
residual-race simulation, is
\begin{align}\label{eq:suppl-S3-prefix-oracle}
 \mathbb E[O(K)\mid C]
 =1-\sum_{j=\max(0,K-m)}^{\min(K,n_0-g_0)}
 \frac{\binom{n_0-g_0}{j}\binom{m}{K-j}}{\binom nK}
 \frac{(A_{\neg})_{K-j}}{(A)_{K-j}},
\end{align}
where $(a)_r=a(a+1)\cdots(a+r-1)$ and $(a)_0=1$.
Indeed a K-subset with $j$ old non-target candidates must also have
$K-j$ future non-target draws; the posterior expectation of their
joint probability is a beta moment
$\mathbb E[(1-\theta_{\rm target})^{K-j}\mid C]
 =(A_{\neg})_{K-j}/(A)_{K-j}$.
This is the correct control-variate anchor for the oracle \emph{and}
for the associated E0/E1 loss; membership of the single residual
winner in the target classes is \emph{not} an oracle.

\paragraph{A sharp logical obstruction from finite moments.}
Let a population of images have unknown mixing law $\mu$ for the
image-specific probability $p\in[0,1]$ of generating raw GOOD.
For $n$ independent runs conditional on image,
\begin{equation}\label{eq:suppl-four-moments}
 \Pr(G_n=j)=\binom nj\int_0^1p^j(1-p)^{n-j}\mu(dp).
\end{equation}
The $n$-run count law determines only the first $n$ moments of
$\mu$. Yet the population oracle at depth $K$ is
$O_K(\mu)=\int[1-(1-p)^K]\mu(dp)$, with limiting ceiling
$\Pr_\mu(p>0)$. The following two mixing laws have \emph{identical}
four-run count probabilities, including in the infinite-image
limit:
\begin{equation}\label{eq:suppl-radau-witness}
 \mu_A=\operatorname{Uniform}[0,1],\qquad
 \mu_B=\frac19\delta_0+
 \frac{16+\sqrt6}{36}\delta_{(6-\sqrt6)/10}
 +\frac{16-\sqrt6}{36}\delta_{(6+\sqrt6)/10}.
\end{equation}
They agree on all polynomials of degree at most four by the
classical three-node left Gauss--Radau quadrature formula.
Consequently their four-run count law is identical
($\Pr(G_4=j)=1/5$ for $j=0,\dots,4$), but
\[
 O_{64}(\mu_A)=64/65\approx .984615,\qquad
 O_{64}(\mu_B)\approx .888889,
\]
a difference of \textbf{9.57 percentage points}, and their limiting
ceilings differ by exactly $1/9$. For DAPS/B22's still shallower
$n=2$ prefix, $\operatorname{Uniform}[0,1]$ and
$\tfrac14\delta_0+\tfrac34\delta_{2/3}$ agree on moments of
degree at most two yet have different large-$K$ ceilings.
Neither pair describes the true image population; nor does
failure in a finite number of runs prove that any particular
image has $p=0$.

\paragraph{Data-specific outer bounds, with the correct target.}
An \emph{exploratory} numerical supplement calculation performed
after FAST64 outcome release used only the original FAST four-run
GOOD-count histograms. Under iid images from a cohort population,
construct five marginal 99\% Clopper--Pearson binomial intervals
$[L_j,U_j]$ for the five four-run count-category probabilities.
By Bonferroni, all five bands hold simultaneously with probability
at least 95\%. Optimize the \emph{population}, not fixed-panel,
$O_{64}(\mu)$ over nonnegative probability measures $\mu$ with
\[
 L_j\le\int\binom4j p^j(1-p)^{4-j}\mu(dp)\le U_j,
 \quad j=0,\ldots,4,\qquad\int\mu(dp)=1.
\]
To certify a conservative continuum outer interval from a finite
mesh of spacing $h=10^{-4}$, assign each continuum mass to its
nearest grid point (distance at most $h/2$). The Bernstein
polynomial count integrands have derivative norm at most four;
the target $1-(1-p)^{64}$ has derivative norm at most 64.
Expanding the grid LP count constraints by $2h$ and its minimum and
maximum target by $32h$ therefore encloses all compatible
continuum measures up to numerical solver tolerance. The resulting
\emph{population} outer bounds are
\begin{equation}\label{eq:suppl-population-bounds}
 O_{64}^{\mathrm{population}}\in
 \begin{cases}
 [0.8696,1.0000],&\text{FFHQ B22 (100 images)},\\
 [0.6113,1.0000],&\text{FFHQ enriched DEV80 (80 images)}.
 \end{cases}
\end{equation}
The confidence interpretation requires iid sampling of images
within each cohort; the deliberately selected/enriched DEV80
cohort does not automatically meet this assumption. These are
\emph{not} 95\% bounds for the registered conditional-prefix
fixed-image realized-panel estimands. They also do not estimate
the intrinsic proportion of unsolvable images precisely.

\paragraph{Registered predictions versus observations: report both halves.}
All reported comparisons below evaluate already committed and
registered forecast artifacts \emph{without refitting}.
Counts across $K$ and uncertainty arms are correlated; coverage
fractions are descriptive across a frozen grid, \emph{not}
independent Bernoulli trials or a calibrated frequentist
significance test.
\begin{itemize}
\item \textbf{FFHQ FAST v1 (64-run panel).}
For E0, race predictive-interval coverage is 6/21, versus S2
15/21. Raw/orbit oracles cover only 6/21 and 4/21 points, and
seven images never generate GOOD across 64 runs, six being
INCONSISTENT in every one. With \emph{realized} class composition
held fixed in a strictly post-outcome diagnostic, the frozen primary
race predicts pooled K64 E0 .9526 versus .9389 observed;
the stricter-prior race gives .9379, S2 .8773. S2's better
unconditional coverage arises partly from cancellation of a
pessimistic rank selector with optimistic generation.
\emph{Dated E1 update (2026-11-03):} on the primary
64-run panel the registered race covers 3/21 E1 intervals
versus 21/21 for S2, yet a \emph{post hoc}
selection-only check with frozen parameters, realized
candidate classes, and realized face-test decisions
overpredicts pooled E1 by only 0.6--0.8 percentage
points for race (while S2 underpredicts by 2.5--4.7
points), and the registered race exceeds its
conditional prediction by 2.9--3.5 points,
identifying the generation forecast as the main
source of its unconditional E1 miss.
\item \textbf{SITCOM B16 v1 and separately registered v3.}
Across its primary $n=16$ panels, v1 race covers all 18/18 E0/E1
selection policy points versus 2/18 for S2; raw/orbit oracles
cover 13/18. At pooled K16 the observed
$(E0,E1,O_{\rm raw},O_{\rm orbit})=(.822,.911,.917,.922)$,
versus race forecast $(.853,.908,.929,.942)$. The FAST-informed
v3 has only 17/36 all-policy coverage, compared with 31/36
for v1, as its generator becomes overly pessimistic on SITCOM.
\item \textbf{DAPS F8 v1 and v3.}
V1 race covers 21/28 overall (E0 7/7, E1 5/7, oracles 9/14),
while S2 covers 22/28 including all 14/14 E0/E1 selection
points. At pooled K8 the observed
$(E0,E1,O_{\rm raw},O_{\rm orbit})=(.928,.939,.939,.944)$
versus race $(.943,.963,.962,.968)$. The original orbit
oracle overforecasts ($.9685$ versus $.944$ observed) and v3
covers only 17/28 overall; DAPS does \emph{not} establish a
race-over-S2 coverage advantage.
\item \textbf{ImageNet FAST D1 (16-run panel).}
In the separately registered v2 E1 addendum, race, strict-prior
race and transferred-tail each cover 6/6 E1 interval checks,
while S2 covers 0/6. Here 6/6 means \emph{three K values times
two different uncertainty constructions evaluated on the same
observations}, not six independent held-out experiments.
At K16 E1 is .650, versus race .603 [.520,.682], but the
orbit oracle is only .750 versus predicted .859 (outside the
frozen band), while raw oracle is .700 versus .683 forecast.
Thus E1 coverage alone does not validate the generation model.
\end{itemize}

\paragraph{Scope of the selection claim.}
These results support a \emph{calibrated, GT-free residual
selection explanation} for particular native solvers/cohorts,
with clear failures of the extrapolated image-generation model.
The conditional-on-realized-label analysis is diagnostic after
unblinding, not prospective verification of the joint forecast.
The exact combinatorial and moment statements are mathematical;
Dirichlet image-rate extrapolation, structural-zero priors,
score surrogates, and CLIP orientation rates are statistical
modeling choices. No new general Plackett--Luce theorem,
Fourier-symmetry theorem, or quadrature inequality is claimed.

% ===== DATED A15 ADDENDUM (earlier S3 text above unchanged) =====
% Theory-owned S3 dated addendum, 2026-10-31. Do not refit forecasts.
\subsection{Outcomes registered after October 29}
\label{app:s3-post-registration-20261031}

\paragraph{Prospective tests 2--5: selection versus generation.}
Under the unchanged model-specific 95\% predictive-interval rule,
the class-sensitive race covers \textbf{57/59} E0/E1 selection
points across tests 2--5, compared with \textbf{17/59}
for exchangeable ranks S2. By test, these counts are ImageNet
FAST with a four-run prefix (9/9 vs 0/9), SITCOM B16
(18/18 vs 2/18), DAPS F8 (12/14 vs 14/14), and
ImageNet FAST with a sixteen-run prefix (18/18 vs 1/18).
Test 1, FFHQ FAST64, is excluded from this four-test tally,
not from the overall reported results: that test's class
generation miss spoiled unconditional race coverage.
These points share images and nested K-subsets, so the coverage tally
is descriptive rather than a binomial test with 59 independent trials.
The DAPS exception (12/14 vs 14/14) precludes claiming
that race wins in every family. We also report
that generation-oracle forecasts were optimistic in
\textbf{four of five} tests. These are distinct,
precommitted observables: successful E0/E1 interval coverage
does not prove that generation and ranking are each calibrated.

\paragraph{Third prospective test: registered v4 and ImageNet XP-D4.}
Before XP-D4 runs 16--63 were opened, v4 was registered
(theory freeze 149331c) after refitting on only
100 fixed ImageNet images' first sixteen runs.
It used the original five-class Dirichlet and four
rank variants, with \emph{no structural-zero spike}.
The primary completed panel has $n_{\max}=32$,
$K=16,24,32$; the secondary has $n_{\max}=64$,
$K=32,48,64$. Race covers \textbf{18/18} E0/E1
points against S2's \textbf{1/18}, but the raw oracle
covers only \textbf{2/6}, and the correct-orbit oracle
covers \textbf{0/6}. At full-panel $K=32$,
the observed orbit oracle is $.750$ versus $.840$
forecast; at $K=64$ it is $.780$ versus $.899$.
There are 25 images with no correct-orbit candidate
among 16 runs and still 22 among 64. The observed
long tail remains an out-of-sample failure of the
registered generator, not a reason to revise v4.
Errors in the orbit oracle and the unconditional
selector--oracle loss may cancel in the well-covered E1.

\paragraph{DAPS F8: an important counterpoint.}
On the held eight-run FFHQ DAPS panel (100 B22 plus
80 DEV80 images), v1 race covers \textbf{21/28}
reported E0/E1/oracle points (E0 7/7, E1 5/7,
raw oracle 5/7, orbit oracle 4/7); S2 covers
\textbf{22/28}, including all 14 E0/E1 points;
the separately frozen FAST-informed v3 covers
\textbf{17/28}. At pooled $K=8$, the observed
$(E0,E1,O_{\rm raw},O_{\rm orbit})$ is
$(.928,.939,.939,.944)$ against v1 race
$(.943,.963,.962,.968)$. The orbit forecast
$.9685$ has a lower 95\% endpoint $.94475$,
slightly above observed $.944$. DAPS has relatively
few historically passing-but-wrong CONSISTENT candidates,
so repaired E1 is less discriminating between
S2 and race than on SITCOM; their E0 predictions
remain different. V3 better centers this one
orbit-oracle value but has worse total coverage.
Neither v3 nor S2 is retrospectively promoted as
the registered primary v1 race.

\paragraph{Sixteen moments still do not identify the rare tail.}
Under the idealized Bernoulli-mixture model, the entire
sixteen-run count law determines only integrals of
polynomials of degree at most sixteen against the
unknown image-success-rate distribution $\mu$.
The population oracle
$O_{64}=\int[1-(1-p)^{64}]\mu(dp)$ instead
requires a degree-64 functional. Distinct
mixing laws can therefore match all sixteen-run
count probabilities yet disagree on this oracle
and on the structural-zero mass, except in
special degenerate moment cases.
For an explicit witness, let $\mu_A$ be uniform on
$[0,1]$ and let $\mu_B$ be its classical nine-node
left Gauss--Radau quadrature (node 0 and the eight
shifted roots of $P_8^{(0,1)}$). The quadrature
is exact through degree sixteen, so both produce
the same sixteen-run Binomial count law; nevertheless
$\mu_B(\{0\})=1/81$ and $\mu_A(\{0\})=0$.
They differ in $O_{64}$ by approximately
$9.39\times10^{-4}$ and in the infinite-run
oracle ceiling by $1/81$. This witness proves
\emph{nonidentification}, not the magnitude of
v4's actual $0.09$--$0.12$ orbit forecast errors.
With only 100 finite images, moment estimation
is additionally noisy; rank/orbit categories
introduce further uncertainty. Even a positive
$p=.002$ gives a sixteen-failure probability
$(1-.002)^{16}\simeq.968$.
A rare near-zero-success component \emph{influences}
low-order moments but is generally \emph{not uniquely
identified} by them. Sixteen observed runs thus
cannot guarantee accurate sixty-four-run coverage.

\paragraph{Registration and scope.}
All outcomes above were released only after their
respective forecast freezes. This dated addendum
preserves the earlier S3 text and does not revise
the Dirichlet model, rank rates, or historical
Clopper--Pearson/LP bounds. The quadrature identity
is classical; none of these observations proves that
a specific real image is mathematically unsolvable.

% ===== DATED A19 TEST-6 ADDENDUM (older S3 above unchanged) =====
% DATED POST-REGISTRATION ADDENDUM: 2026-11-02.
% Theory-owned; do not edit any frozen v5 forecast or its comparison rule.
\subsection{Test 6: first prospective test on image-disjoint ImageNet images}
\label{app:s3-test6-xpn-20261102}

\paragraph{Registration and scoring.}
The v5 code was frozen before the new ImageNet images' outcomes
were available (code commit \texttt{81003cc}), and the numerical
v5 forecast was sealed and pushed on 2026-11-01 at 19:53:58 UTC
(XP commit \texttt{55ba86b}), before any run with index
$4$--$63$ was scored. The model was fitted only on runs
$0$--$3$ of the 100 \emph{image-disjoint} test images.
The registered primary comparison is race-model predictive-interval
coverage for E0 and E1 (class-Beta decision model), across
$K\in\{4,8,16,32,64\}$ from the \emph{complete}
64-run panel. Neither its intervals nor its model-selection
rule is revised in light of the results.

On the ten predeclared selection points, race covers
\textbf{10/10}, followed by the stronger-regularization
race variant \textbf{7/10}, the transferred tail variant
\textbf{5/10}, and exchangeable ranks S2 \textbf{0/10}.
Their mean absolute standardized errors are
$1.236$, $1.329$, $2.258$, and $7.441$, respectively.
Race is the winner under the frozen rule.
The five $K$-points share the \emph{same} images, candidate panels,
and nested subset evaluations; the ten coverage indicators are
correlated, not ten independent Bernoulli trials.
Combining the prospectively scored tests 2--6 gives
race \textbf{67/69} versus S2 \textbf{17/69}
selection-coverage points; this pooled descriptive tally
is likewise correlated and does not establish exact calibration
of a standalone generative mechanism.

\paragraph{Observed reconstruction and oracle curves.}
The complete-panel successes are
\[
\begin{array}{c|ccccc}
K&4&8&16&32&64\\ \hline
E0&.359&.415&.463&.501&.520\\
E1&.456&.498&.542&.585&.610\\
O_{\rm raw}&.470&.564&.634&.690&.730\\
O_{\rm orbit}&.569&.624&.669&.711&.750
\end{array}
\]
For race at $K=64$, the frozen E0 forecast is
$.462\ [ .378,.544]$ versus $.520$ observed,
and E1 is $.626\ [ .534,.711]$ versus $.610$.
By contrast, the correct-orbit oracle forecast
$.880\ [ .821,.931]$ misses the observed $.750$,
and the registered orbit-oracle interval misses at
\emph{all five} $K$-values (\textbf{0/5}).
The raw-GOOD oracle covers \textbf{3/5} values and
is \emph{underpredicted} at $K=4,8$:
$.417$ versus $.470$ and $.510$ versus $.564$.
The orbit-generation miss is therefore not a simple
uniform downward shift in all oracle success rates.
The correct-orbit oracle was optimistic in five
of the six prospectively reported tests.

\paragraph{Which image-level dependency was missed?}
Let $G_K$ mean that a uniform $K$-subset includes at
least one raw-GOOD run, and let $O_K$ mean that it includes
a correct-orbit but non-GOOD run. Then, exactly,
\begin{equation}\label{eq:s3-test6-orbit-only}
 O_{\rm orbit}(K)-O_{\rm raw}(K)
 =\Pr(O_K\cap G_K^c).
\end{equation}
At $K=64$, the registered v5 difference is
$.164$, but its observed value is \emph{only}
$.750-.730=.020$. Thus the model expected about
$16.4$ of 100 images to be rescued \emph{only} by
an orientation-orbit candidate, compared with
two actually observed. Independent inspection
of the authorized 64-run class strings gives
\textbf{67} images with both a GOOD and a non-GOOD
correct-orbit run, \textbf{6} with GOOD only,
\textbf{2} with orbit-only, and \textbf{25} with
neither. Across the 100 images, the respective
64-run GOOD and non-GOOD correct-orbit class counts
have descriptive Pearson correlation approximately
$+0.369$. In particular, $67/69$ images containing
an orbit-error run also contain a raw-GOOD run.
The low increment in
\eqref{eq:s3-test6-orbit-only} therefore reflects
\emph{co-occurrence across images}, not the absence
of orbit-error candidates.

The v5 hierarchical five-class Dirichlet model
and its observed-prefix conditioning did not
represent this image-level overlap adequately
when extrapolated to 64 runs. A single Dirichlet
distribution for a new image's class probabilities
also imposes negative covariance between its
distinct class-probability coordinates, whereas
the observed long-run class counts are positively
associated across images. This is evidence of
a shared easy/hard propensity or other dependence
outside that simple prior, although the observed
count correlation alone does not identify a unique
alternative joint mixing law or causal mechanism.
It is a \emph{model-misspecification diagnosis}:
the primary selector's accurate coverage can coexist
with offsetting errors in the oracle and
oracle-to-selector selection-loss models.

\paragraph{Relation to finite-moment non-identifiability.}
If a fixed image has per-run class probabilities
$\theta_G$ (raw GOOD) and
$\theta_O$ (other correct-orbit), then the
idealized population increment is the joint
mixing functional
\begin{equation}\label{eq:s3-test6-joint-functional}
 D_K=\mathbb E\!\left[
 (1-\theta_G)^K-
 (1-\theta_G-\theta_O)^K
 \right].
\end{equation}
Even the \emph{complete distribution} of
four-run five-class counts determines only
polynomial moments of total degree at most four
of the joint law of $(\theta_G,\theta_O,\ldots)$,
not the degree-$64$ functional
\eqref{eq:s3-test6-joint-functional}.
The C-T8 truncated-moment obstruction is
therefore a reason why correct long-run
coverage cannot be guaranteed from short-run
counts, and applies more broadly to the
\emph{joint} orbit-only event.
However, it neither predicts the sign nor
explains the magnitude of v5's $14.4$-point
orbit-only excess. The concrete, independently
released 64-run overlap table demonstrates
that the \emph{particular registered Dirichlet
extrapolation} fails on the true new-image
joint dependence. Non-identifiability is a
limitation of available calibration information;
misspecification describes this model's
actual out-of-sample prediction error.
Neither statement implies mathematical
unsolvability for an individual image:
$25/100$ images simply have no correct-orbit
candidate in the observed 64-run panel.

\paragraph{Separate stopping pre-score (not v5 fitting).}
The independently frozen stopping prediction
called N-H1a a pass at $N_{\max}=16$
(central gain $+0.2$ points, stated range
$[-1.8,+2.2]$), but the observed
equal-mean-runs gain was $-1.516$ points,
so the \emph{binary prospective call failed}
despite the observation lying inside its
informal range.
N-H1b at $N_{\max}=64$ was correctly called
to pass its negative-gain criterion
(predicted $-2.7$, observed $-3.102$ points).
Image-bootstrap intervals are
$[-3.60,+0.47]$ and $[-5.96,-0.10]$
points, respectively, on the same 100 images.
These stopping observations are distinct from
the v5 E0/E1 predictive-interval scoring,
and are not used to refit any forecast.
